\section{Mathematical Formulation} \label{sec:math}

The FSA-v2 model is defined as a system of three coupled Itô Stochastic Differential Equations (SDEs).

\subsection{State Space and Drift Coefficients}

The state vector is $y(t) = (B(t), F(t), A(t))^T$. The dynamics are governed by the following equations:

\begin{align}
  \dd B &= \left[\,\kappa_B \cdot \frac{1+\varepsilon_A A}{1+\varepsilon_A A_{\mathrm{typ}}}\,\Phi(t)
                  - \frac{B}{\tau_B}\,\right]\dd t
            + \sigma_B \sqrt{B(1-B)}\,\dd W_B \\
  \dd F &= \left[\,\kappa_F\,\Phi(t)
                  - \frac{1+\lambda_A A}{1+\lambda_A A_{\mathrm{typ}}}\,
                    \frac{F}{\tau_F}\,\right]\dd t
            + \sigma_F \sqrt{F}\,\dd W_F \\
  \dd A &= \bigl[\,\mu(B,F)\,A - \eta\,A^3\,\bigr]\,\dd t
            + \sigma_A \sqrt{A}\,\dd W_A
\end{align}

where $\mu(B,F)$ is the Stuart-Landau bifurcation parameter:
\begin{equation}
  \mu(B,F) = \mu_0 + \mu_B B - \mu_F F - \mu_{FF}\,(F - F_{\mathrm{typ}})^2
\end{equation}

\subsection{Diffusion and Boundary Conditions}

The model employs state-dependent noise to maintain physiological boundaries:
\begin{itemize}
    \item \textbf{Jacobi Diffusion for B:} The term $\sigma_B \sqrt{B(1-B)}$ ensures $B$ stays within the unit interval $[0, 1]$. The diffusion vanishes as $B \to 0$ or $B \to 1$.
    \item \textbf{CIR Diffusion for F and A:} The terms $\sigma_F \sqrt{F}$ and $\sigma_A \sqrt{A}$ ensure $F, A \ge 0$. The diffusion vanishes as the state approaches zero, preventing negative values (Cox-Ingersoll-Ross process).
\end{itemize}

\subsection{The G1 Reparametrization}

The G1 form redefines parameters to ensure structural identifiability. By scaling the couplings by their values at the operating point $(A_{\mathrm{typ}}, F_{\mathrm{typ}})$, we create a system where:
\begin{enumerate}
    \item $\kappa_B$ is the \textit{effective} fitness gain at $A_{\mathrm{typ}}$.
    \item $\tau_F$ is the \textit{effective} fatigue time constant at $A_{\mathrm{typ}}$.
    \item $\mu_F$ is the \textit{local slope} of the growth rate with respect to fatigue at $F_{\mathrm{typ}}$.
\end{enumerate}
This rotation ensures that the primary parameters are strongly identified by the data, while the residuals ($\varepsilon_A, \lambda_A, \mu_{FF}$) capture higher-order non-linearities.

\subsection{Observation Model}

The latent states are projected onto four observable channels:
\begin{align*}
  \mathrm{HR}(t) &\sim \mathcal{N}(\text{HR}_{\text{base}} - \kappa_B^{HR}B + \alpha_A^{HR}A + \beta_C^{HR}C(t), \sigma_{HR}^2) \\
  P(\text{sleep}(t)=1) &= \text{sigmoid}(k_C C(t) + k_A A - \tilde{c}) \\
  S(t) &\sim \mathcal{N}(S_{\text{base}} + k_F F - k_A^S A + \beta_C^S C(t), \sigma_S^2) \\
  \log(\text{steps}(t)+1) &\sim \mathcal{N}(\mu_0^{\text{step}} + \beta_B^{\text{step}}B - \beta_F^{\text{step}}F + \beta_A^{\text{step}}A + \beta_C^{\text{step}}C(t), \sigma_{\text{st}}^2)
\end{align*}
where $C(t)$ is a fixed 24-hour circadian cycle.
